% !TeX root = ../../Book3.tex
\section{Schwarz–Pick 引理}
\label{b3:sec:0603}

Schwarz 引理处理固定原点的圆盘映射。要比较任意两点，先把一个源点及其像分别送到原点，再对复合函数使用同一引理。撤销坐标后，普通的模估计就变成伪双曲距离压缩，原点导数估计则变成随位置变化的导数界。这两个结论既能判断插值数据是否相容，也能控制有内部不动点的迭代。

\subsection{Schwarz–Pick 估计与两点插值}
\label{b3:subsec:060301}

以下仍记 \(\phi_a(z)=(z-a)/(1-\overline a z)\)，其中 \(a\in\mathbb D\)。前节定义的距离为 \(\delta(z,w)=|\phi_w(z)|\)。\term[schwarz-pick-yinli]{Schwarz--Pick 引理}[Schwarz--Pick lemma]说明，全纯圆盘映射虽然未必压缩欧氏距离，却总压缩这个距离。

\Needspace{6\baselineskip}
\begin{theorem}[Schwarz--Pick]\label{b3:thm:06-schwarz-pick}
设 \(f:\mathbb D\rightarrow\mathbb D\) 全纯。则对所有 \(z,w\in\mathbb D\)，有
\[
 \delta(f(z),f(w))\le\delta(z,w),
 \quad
 \frac{|f'(w)|}{1-|f(w)|^2}\le\frac1{1-|w|^2}.
\]
若在某一点的导数不等式中取等，或在某一对不同点的距离不等式中取等，则 \(f\) 是圆盘自同构。反过来，圆盘自同构在上述所有不等式中都取等。
\end{theorem}
\begin{proof}
固定 \(w\in\mathbb D\)，令
\[
 F=\phi_{f(w)}\circ f\circ\phi_w^{-1}.
\]
两端的分式映射都是圆盘自同构，且 \(\phi_w^{-1}(0)=w\)，故 \(F:\mathbb D\rightarrow\mathbb D\) 全纯并满足 \(F(0)=0\)。Schwarz 引理给出
\[
 |F(\phi_w(z))|\le|\phi_w(z)|.
\]
左端是 \(|\phi_{f(w)}(f(z))|\)，所以这就是距离不等式。

对导数，圆盘自同构的公式给
\[
 \phi_{f(w)}'(f(w))=\frac1{1-|f(w)|^2},
 \quad (\phi_w^{-1})'(0)=1-|w|^2.
\]
于是链式法则与 Schwarz 的原点导数估计给出
\[
 |F'(0)|=
 \frac{|f'(w)|\cdot(1-|w|^2)}{1-|f(w)|^2}\le1.
\]
分母均为正，整理即得所需导数式。

若导数式取等，则 \(|F'(0)|=1\)，Schwarz 等号情形使 \(F\) 为旋转。若距离式对 \(z\ne w\) 取等，则 \(\phi_w(z)\ne0\)，同一等号情形仍使 \(F\) 为旋转。撤销两个自同构，\(f\) 便是圆盘自同构。反方向可分别对 \(f\) 和 \(f^{-1}\) 使用已证不等式，两个相反方向的估计迫使距离和导数处处取等。
\end{proof}

“不同点”的条件使距离等号具有内容；取同一点时，任何函数都会给出 \(0=0\)。常值函数也满足两个压缩式，但不在导数式或不同点的距离式中取等。若起初只有 \(|f|\le1\)，则先用最大模原理区分：\(f\) 要么是模一常值，要么处处落在开圆盘中。模一常值的分母 \(1-|f|^2\) 为零，须直接处理。

\ProseParagraphBreak

固定 \(w\)，把距离式除以 \(|z-w|\)，再令 \(z\rightarrow w\)，也会得到导数式。其分母分别趋于 \(1-|f(w)|^2\) 与 \(1-|w|^2\)。不过，一处导数界本身不会直接给出任意两点的比较；上面的坐标归一化同时控制了整个圆盘中的点值。

\Needspace{6\baselineskip}
\begin{proposition}[固定函数值的导数界]\label{b3:prop:06-value-derivative-bound}
设 \(a\in\mathbb C\)、\(R,M>0\)，\(f:D(a,R)\rightarrow D(0,M)\) 全纯，并记 \(c=f(a)\)。则
\[
 |f'(a)|\le\frac{M^2-|c|^2}{MR}.
\]
对每个 \(c\in D(0,M)\)，都存在具有该中心值并达到导数上界的映射。
\end{proposition}
\begin{proof}
在\eqref{b3:eq:06-center-derivative}中令目标圆心为零、目标半径为 \(M\)，便得所列导数界。为了确定它可以达到，任取 \(|\eta|=1\)，令
\[
 f(z)=M\cdot\phi_{c/M}^{-1}\left(\eta\frac{z-a}{R}\right).
\]
这个映射双全纯到 \(D(0,M)\)，在 \(a\) 取值 \(c\)，导数为 \(\eta(M^2-|c|^2)/(MR)\)，故达到上界。
\end{proof}

例如 \(M=R=1\)、\(a=0\)、\(c=1/2\) 时，导数模至多为 \(3/4\)，而不利用中心值的 Cauchy 估计只给 \(1\)。映射 \((z+1/2)/(1+z/2)\) 在原点的值为 \(1/2\)，导数为 \(3/4\)，说明这个改进是准确的。

\ProseParagraphBreak

两点插值把估计的必要条件变成存在性的完整判据。构造时先把第一组插值数据移到零点，余下只需用一个线性映射把第二个源值送到第二个目标值。

\Needspace{6\baselineskip}
\begin{proposition}[两点插值]\label{b3:prop:06-two-point-interpolation}
给定不同的 \(z_1,z_2\in\mathbb D\) 及 \(w_1,w_2\in\mathbb D\)。存在全纯映射 \(f:\mathbb D\rightarrow\mathbb D\)，使 \(f(z_j)=w_j\)、\(j=1,2\)，当且仅当
\[
 \delta(w_1,w_2)\le\delta(z_1,z_2).
\]
若取等，插值映射唯一，且为圆盘自同构；若严格不等，则有无穷多个不同的插值映射。
\end{proposition}
\begin{proof}
必要性来自 Schwarz--Pick。为证明充分性，令
\[
 p=\phi_{z_1}(z_2),\quad q=\phi_{w_1}(w_2),\quad c=q/p.
\]
源点互异使 \(p\ne0\)，条件又给 \(|c|\le1\)。于是
\[
 f(z)=\phi_{w_1}^{-1}\bigl(c\cdot\phi_{z_1}(z)\bigr)
\]
全纯地将圆盘映入自身，并满足两项插值。当 \(|c|=1\) 时，它是自同构。任何其他解的归一化函数 \(F=\phi_{w_1}\circ f\circ\phi_{z_1}^{-1}\) 满足 \(F(0)=0\)、\(F(p)=q\)。若 \(|q|=|p|\)，Schwarz 等号迫使 \(F(\zeta)=c\zeta\)，所以这个解唯一。

若 \(|c|<1\)，对每个满足
\[
 |t|<\frac{1-|c|}{1+|p|}
\]
的复数 \(t\)，定义 \(H_t(\zeta)=c+t(\zeta-p)\)。对 \(\zeta\in\mathbb D\)，有 \(|H_t(\zeta)|\le|c|+|t|\cdot(1+|p|)<1\)，而 \(H_t(p)=c\)。因此 \(F_t(\zeta)=\zeta\cdot H_t(\zeta)\) 都满足 \(F_t(0)=0\)、\(F_t(p)=q\)。撤销坐标得到插值映射 \(f_t\)。不同的 \(t\) 给出不同的多项式 \(F_t\)，故也给出不同的 \(f_t\)。
\end{proof}

取 \(z_1=w_1=0\)、\(z_2=1/2\)。若要求 \(w_2=3/4\)，目标距离大于源距离，无解；若改为 \(w_2=1/4\)，则 \(f(z)=z/2\) 是一个解；若取 \(w_2=1/2\)，唯一解是恒等映射。对三个点，原始数据的成对距离条件还不够，下一小节将用一次除零点步骤找出剩余的相容条件。

\ProseParagraphBreak

严格压缩不意味着全圆盘共有一个小于 \(1\) 的压缩系数。对 \(f(z)=z^2\)，其归一化导数比为
\[
 \frac{|f'(z)|\cdot(1-|z|^2)}{1-|f(z)|^2}
 =\frac{2|z|}{1+|z|^2}<1.
\]
每个内点都严格小于 \(1\)，但沿 \(|z|\uparrow1\) 趋于 \(1\)。只有固定 \(|z|\le r<1\) 后，才能用 \(2r/(1+r^2)<1\) 作共同上界。

\subsection{不动点迭代与插值构造}
\label{b3:subsec:060302}

圆盘映射若有内部不动点，就可以通过共轭把它移到原点。此时 Schwarz 引理使每条轨道留在一个内部紧盘里，紧盘上的严格性足以证明整条迭代轨道收敛。

\Needspace{6\baselineskip}
\begin{proposition}[内部不动点的刚性]\label{b3:prop:06-two-fixed-points}
设 \(f:\mathbb D\rightarrow\mathbb D\) 全纯。若 \(f\) 固定两个不同的内部点，则 \(f\) 为恒等映射。若 \(a\in\mathbb D\)、\(f(a)=a\) 且 \(f'(a)=1\)，则 \(f\) 也为恒等映射。
\end{proposition}
\begin{proof}
若有两个不同不动点 \(a,b\)，距离不等式在该点对取等，故 \(f\) 是自同构。映射 \(g=\phi_a\circ f\circ\phi_a^{-1}\) 固定零点，所以是旋转 \(g(\zeta)=\eta\zeta\)。它还固定非零点 \(\phi_a(b)\)，故 \(\eta=1\)，撤销共轭即得 \(f=\mathrm{id}\)。

对第二项，同样的共轭满足 \(g(0)=0\)，并由链式法则得到 \(g'(0)=f'(a)=1\)。Schwarz 导数等号给出 \(g(\zeta)=\zeta\)，于是 \(f\) 仍是恒等映射。
\end{proof}

\Needspace{6\baselineskip}
\begin{proposition}[内点不动点处的迭代]\label{b3:prop:06-interior-iteration}
设 \(f:\mathbb D\rightarrow\mathbb D\) 全纯，\(a\in\mathbb D\)，且 \(f(a)=a\)。若 \(f\) 不是圆盘自同构，则迭代 \(f^{\circ n}\) 在 \(\mathbb D\) 上局部一致趋于常值 \(a\)。在已有内部不动点的条件下，\(f\) 不是圆盘自同构等价于 \(|f'(a)|<1\)。
\end{proposition}
\begin{proof}
令 \(g=\phi_a\circ f\circ\phi_a^{-1}\)。它固定零点，且 \(g'(0)=f'(a)\)。Schwarz 引理给 \(|g'(0)|\le1\)，取等当且仅当 \(g\) 为旋转；这也恰是 \(f\) 为圆盘自同构的情形。因此题设等价于 \(|f'(a)|<1\)。

在原点补上导数后，函数 \(h(\zeta)=g(\zeta)/\zeta\) 全纯且 \(|h|\le1\)。若它在某个内点取模 \(1\)，最大模原理使它为单位常数，\(g\) 就成为旋转，与假设矛盾。所以对每个 \(0<r<1\)，都有
\[
 q_r=\max_{|\zeta|\le r}|h(\zeta)|<1.
\]
Schwarz 引理使 \(g\) 把 \(\overline{D(0,r)}\) 映入自身，故可反复使用同一个常数，归纳得到
\[
 |g^{\circ n}(\zeta)|\le q_r^n\cdot|\zeta|
 \quad(|\zeta|\le r).
\]
对原圆盘中的固定紧集 \(K\)，选 \(r<1\) 使 \(\phi_a(K)\subseteq\overline{D(0,r)}\)。因为
\[
 \phi_a^{-1}(u)-a=\frac{(1-|a|^2)u}{1+\overline a u},
\]
共轭关系给出
\[
 \sup_{z\in K}|f^{\circ n}(z)-a|
 \le\frac{1-|a|^2}{1-|a|\cdot r}\,q_r^n r
 \longrightarrow0.
\]
分母为正，且不随 \(n\) 改变，故得到紧集上的一致收敛。
\end{proof}

旋转是这一结论的准确例外。若共轭后的映射为 \(g(\zeta)=\eta\zeta\)、\(|\eta|=1\)，则迭代为 \(\eta^n\zeta\)。当 \(\eta\ne1\) 时，任何非零轨道都不收敛：若极限为 \(L\)，其模非零，而相邻项关系又给 \(L=\eta L\)，矛盾。单位根情形的轨道是周期的；\(\eta=1\) 时则每个点固定。没有内部不动点的映射可能有趋近边界的轨道，那属于后续 Denjoy--Wolff 理论，本处不据此作结论。

\ProseParagraphBreak

例如 \(g(z)=z^2\) 的第 \(n\) 次迭代为 \(z^{2^n}\)，在 \(|z|\le r<1\) 上误差至多 \(r^{2^n}\)，却不在整个开圆盘上一致趋于零，其模的全域上确界始终为 \(1\)。固定紧集使所有迭代点共享一个严格小于 \(1\) 的估计。

\ProseParagraphBreak

把归一化后的零点除去，也能逐步处理插值数据。一次取值条件被消去以后，剩余自由度由另一个模不超过 \(1\) 的全纯函数表示。这称为\term[schur-buzhou]{Schur 步骤}[Schur step]。

\Needspace{6\baselineskip}
\begin{proposition}[Schur 步骤]\label{b3:prop:06-schur-step}
给定 \(c\in\mathbb D\)，令 \(\phi_c(w)=(w-c)/(1-\overline c w)\)。若 \(f:\mathbb D\rightarrow\mathbb D\) 全纯且 \(f(0)=c\)，则
\[
 g(z)=\frac{\phi_c(f(z))}{z}\quad(z\ne0)
\]
在原点可去，补定义后全纯且 \(|g(z)|\le1\)，并有
\[
 g(0)=\frac{f'(0)}{1-|c|^2}.
\]
反过来，对每个在 \(\mathbb D\) 上全纯且模不超过 \(1\) 的函数 \(g\)，公式
\[
 f(z)=\phi_c^{-1}\bigl(z\cdot g(z)\bigr)
     =\frac{c+z\cdot g(z)}{1+\overline c z\cdot g(z)}
\]
定义一个满足 \(f(0)=c\) 的全纯圆盘映射。这两个操作互为逆操作。
\end{proposition}
\begin{proof}
函数 \(\phi_c\circ f\) 固定零点，Schwarz 引理给 \(|\phi_c(f(z))|\le|z|\)。其 Taylor 展开的一次系数为 \(\phi_c'(c)\cdot f'(0)=f'(0)/(1-|c|^2)\)，所以除以 \(z\) 后可去，且原点补值如上。模估计在非零点成立，在原点由连续性成立。

反过来，\(|z\cdot g(z)|\le|z|<1\)，因此即使 \(g\) 是单位常数，乘积也始终落在开圆盘内。分母的模至少为 \(1-|c|\cdot|z|>0\)，复合逆自同构合法，给出全纯圆盘映射及 \(f(0)=c\)。两边施加 \(\phi_c\) 得到 \(\phi_c(f(z))=z\cdot g(z)\)，从而两个操作互逆。
\end{proof}

如果 \(|f'(0)|=1-|c|^2\)，则 \(|g(0)|=1\)，最大模原理迫使 \(g\) 是单位常数，原函数因而是自同构。若 \(|g(0)|<1\)，则 \(g\) 全部取值于开盘，才可继续作下一步。例如 \(f(0)=f'(0)=0\) 时，连续两次 Schwarz 除零点得到 \(|f(z)|\le|z|^2\)；若再有 \(f''(0)=2\eta\)、\(|\eta|=1\)，第二个商被迫恒为 \(\eta\)，所以 \(f(z)=\eta z^2\)。

\Needspace{6\baselineskip}
\begin{proposition}[二阶 Taylor 数据的相容性]\label{b3:prop:06-second-jet}
给定 \(c,d,e\in\mathbb C\)、\(|c|<1\)，记
\[
 A=1-|c|^2,\quad p=\frac dA,\quad
 q=\frac eA+\frac{\overline c d^2}{A^2}.
\]
存在全纯映射 \(f:\mathbb D\rightarrow\mathbb D\)，使其原点展开为
\[
 f(z)=c+dz+ez^2+O(z^3),
\]
当且仅当以下两种情形之一成立：\(|p|<1\) 且 \(|q|\le1-|p|^2\)；或者 \(|p|=1\) 且 \(q=0\)。这里 \(e=f''(0)/2\)。
\end{proposition}
\begin{proof}
先作 Schur 步骤，令 \(g=\phi_c(f)/z\)。为了辨认 \(g\) 的前两项，记 \(t=f-c=dz+ez^2+O(z^3)\)。因为 \(A>0\)，原点附近有
\[
 \begin{aligned}
 \phi_c(f)&=\frac{t}{A-\overline c t}
   =\frac tA+\frac{\overline c}{A^2}t^2+O(z^3)\\
 &=\frac dA z+
   \left(\frac eA+\frac{\overline c d^2}{A^2}\right)z^2+O(z^3).
 \end{aligned}
\]
所以 \(g(z)=p+qz+O(z^2)\)。Schur 步骤给 \(|g|\le1\)，先有 \(|p|\le1\)。若 \(|p|=1\)，最大模原理使 \(g\equiv p\)，故 \(q=0\)。若 \(|p|<1\)，最大模原理排除其他点取模 \(1\)，于是 \(g\) 是开圆盘自映射，Schwarz--Pick 在原点给出 \(|q|\le1-|p|^2\)。这证明了必要性。

为证明充分性，若 \(|p|<1\)，令
\[
 k=\frac q{1-|p|^2},\quad
 g(z)=\phi_p^{-1}(kz)=\frac{p+kz}{1+\overline p kz}.
\]
条件保证 \(|k|\le1\)，所以 \(g\) 把开圆盘映入开圆盘，且 \(g(0)=p\)、\(g'(0)=(1-|p|^2)k=q\)。若 \(|p|=1\)、\(q=0\)，则取 \(g\equiv p\)。两种情形都可以令 \(f=\phi_c^{-1}(zg)\)，Schur 逆步骤保证 \(f:\mathbb D\rightarrow\mathbb D\) 全纯。

还须核对恢复的二阶项。逆分式在零点附近展开为
\[
 \phi_c^{-1}(u)=c+A u-A\overline c u^2+O(u^3).
\]
代入 \(u=z\cdot g(z)=pz+qz^2+O(z^3)\)，得到
\[
 f(z)=c+A pz+A(q-\overline c p^2)z^2+O(z^3).
\]
定义给出 \(Ap=d\)、\(A(q-\overline c p^2)=e\)，因此所构造函数具有全部指定数据。
\end{proof}

当 \(c=0\) 时，条件变成 \(|d|\le1\)、\(|e|\le1-|d|^2\)。例如 \(d=3/4\)、\(e=1/2\) 不相容，因为 \(1/2>7/16\)，尽管两个系数单独都满足 Cauchy 的模上界 \(1\)。当 \(c=1/2\)、\(d=3/4\) 时，第一步已经给 \(p=1\)，于是 \(q=0\) 强制 \(e=-3/8\)。这正是 \((z+1/2)/(1+z/2)\) 的二次系数，显示分式归一化中的 \(\overline c d^2\) 项不可丢掉。

\ProseParagraphBreak

三点插值也能通过一次 Schur 步骤变成两点问题。除零点后出现的指定值可能位于单位圆上，因此需要先处理这个边界情形，再使用开圆盘中的两点判据。

\Needspace{6\baselineskip}
\begin{proposition}[三点插值的 Schur 归约]\label{b3:prop:06-three-point-reduction}
设 \(a_1,a_2,a_3\in\mathbb D\) 两两不同，\(b_1,b_2,b_3\in\mathbb D\)。记 \(\phi_a(z)=(z-a)/(1-\overline a z)\)，并令
\[
 p_j=\phi_{a_1}(a_j),\quad
 q_j=\frac{\phi_{b_1}(b_j)}{p_j}\quad(j=2,3).
\]
存在全纯映射 \(f:\mathbb D\rightarrow\mathbb D\)，使 \(f(a_j)=b_j\)、\(j=1,2,3\)，当且仅当以下之一成立：\(q_2,q_3\in\mathbb D\) 且 \(\delta(q_2,q_3)\le\delta(p_2,p_3)\)；或者 \(q_2=q_3\) 且 \(|q_2|=1\)。
\end{proposition}
\begin{proof}
源点互异使 \(p_2,p_3\) 非零且不同，故各商有定义。若插值映射存在，则
\[
 G=\phi_{b_1}\circ f\circ\phi_{a_1}^{-1}
\]
是固定零点的圆盘映射。Schur 步骤使 \(H(\zeta)=G(\zeta)/\zeta\) 在零点可去，且 \(|H|\le1\)。插值条件等价于 \(H(p_j)=q_j\)、\(j=2,3\)。如果某个 \(|q_j|>1\)，便不可能；如果某个 \(|q_j|=1\)，最大模原理迫使 \(H\) 恒为这个单位数，所以另一个 \(q_j\) 必须相同。剩下 \(q_2,q_3\in\mathbb D\) 的情形，\(H\) 也处处在开盘内，两点插值判据给出所列距离条件。

反过来，若 \(q_2=q_3\) 为单位数，取 \(H\equiv q_2\)。若两值都在开盘且满足距离条件，可以明确取
\[
 k=\frac{\phi_{q_2}(q_3)}{\phi_{p_2}(p_3)},
 \quad H(\zeta)=\phi_{q_2}^{-1}\bigl(k\cdot\phi_{p_2}(\zeta)\bigr).
\]
因为 \(p_2\ne p_3\)，分母非零，而条件给 \(|k|\le1\)。这就是两点命题中的构造，满足 \(H(p_j)=q_j\)。两种情形均可定义
\[
 f(z)=\phi_{b_1}^{-1}\Bigl(
       \phi_{a_1}(z)\cdot H\bigl(\phi_{a_1}(z)\bigr)\Bigr).
\]
乘积的模严格小于 \(1\)，包括 \(H\) 为单位常数的情形，故所得函数全纯且取值于开盘。在 \(a_1\) 处乘积为零；在另外两点处，它等于 \(p_jq_j=\phi_{b_1}(b_j)\)。代入逆分式就恢复全部三项插值。
\end{proof}

若约化函数被迫为单位常数，或最后两点的距离不等式取等，约化后的 \(H\) 唯一，原映射也唯一。若最后两点严格不等，前面的构造给出无穷多个不同的 \(H\)，逆步骤也给出不同的原映射。

\ProseParagraphBreak

例如要求 \(f(0)=0\)、\(f(1/2)=f(-1/2)=2/5\)。每一对原始数据都满足 Schwarz--Pick：中心到另外两点的目标距离为 \(2/5<1/2\)，后两个目标之间的距离为零。但除去已知零点后，\(H=f/z\) 必须在 \(1/2,-1/2\) 分别取值 \(4/5,-4/5\)。这不可能，因为
\[
 \delta(4/5,-4/5)=\frac{40}{41}
 >\frac45=\delta(1/2,-1/2).
\]
共同的 Schur 归约揭示了各点数据之间的额外约束，不能把不同点对各自存在的解直接拼接。

\subsection{等号刚性与正实部估计}
\label{b3:subsec:060303}

内部的精确等号已经迫使映射成为自同构；若导数只是接近极值，还可以量化函数在紧集上离极端映射有多远。设 \(f:\mathbb D\rightarrow\mathbb D\) 全纯，\(f(0)=0\)，且 \(f'(0)=a\in[0,1]\)。当 \(a<1\) 时，商 \(h=f/z\) 补定义后满足 \(h(0)=a\)、\(|h|<1\)。Schwarz--Pick 给 \(\delta(h(z),a)\le|z|\)。固定 \(0<r<1\)，对 \(|z|\le r\) 有
\[
 \begin{aligned}
 |h(z)-a|&\le r\cdot|1-ah(z)|\\
 &\le r\bigl(1-a^2+a\cdot|h(z)-a|\bigr).
 \end{aligned}
\]
因为 \(1-ar>0\)，移项给出 \(|h-a|\le r(1-a^2)/(1-ar)\)。再由 \(|f(z)-z|\le|z|\cdot(|h(z)-a|+1-a)\)，得到
\[
 \sup_{|z|\le r}|f(z)-z|
 \le\frac{r(1-a)(1+r)}{1-ar}.
\]
若 \(a=1\)，Schwarz 等号直接给 \(f=z\)，同一上界仍成立。因此归一化导数趋于 \(1\) 时，函数在每个固定紧盘上一致趋于恒等。对非零复导数，可先旋转函数使导数正实，再与相应旋转映射比较。

\ProseParagraphBreak

把右半平面映到圆盘，还能将正实部条件转成模约束。这个转换保留的信息比一般的模有界条件更强，因为实部的正性同时控制了圆周平均值和各阶系数。

\Needspace{6\baselineskip}
\begin{proposition}[正实部函数的增长估计]\label{b3:prop:06-positive-real-part}
设 \(p:\mathbb D\rightarrow\mathbb C\) 全纯，\(\operatorname{Re}p(z)>0\) 对全部 \(z\in\mathbb D\) 成立，且 \(p(0)=1\)。则对每个 \(z\in\mathbb D\)，记 \(r=|z|\)，有
\[
 \frac{1-r}{1+r}\le\operatorname{Re}p(z)
 \le\frac{1+r}{1-r},
 \quad |p(z)|\le\frac{1+r}{1-r}.
\]
\end{proposition}
\begin{proof}
正实部使 \(p+1\) 无零，且
\[
 |p+1|^2-|p-1|^2=4\operatorname{Re}p>0.
\]
所以 \(q=(p-1)/(p+1)\) 是圆盘映射，\(q(0)=0\)，Schwarz 引理给 \(|q(z)|\le r\)。反解 \(p=(1+q)/(1-q)\)，得到
\[
 \operatorname{Re}p(z)=\frac{1-|q(z)|^2}{|1-q(z)|^2}.
\]
令 \(s=|q(z)|\le r\)。由 \(1-s\le|1-q(z)|\le1+s\)，有
\[
 \frac{1-s}{1+s}\le\operatorname{Re}p(z)\le\frac{1+s}{1-s},
 \quad |p(z)|\le\frac{1+s}{1-s}.
\]
左侧函数随 \(s\) 递减，右侧函数随 \(s\) 递增，代入 \(s\le r\) 即得全部估计。
\end{proof}

函数 \(p_\eta(z)=(1+\eta z)/(1-\eta z)\)、\(|\eta|=1\)，在满足 \(\eta z=r\) 的点达到上界，在满足 \(\eta z=-r\) 的点达到实部下界。因此两个方向的界都可达到。这里的推导只对全纯函数的实部进行，后面的调和函数理论将解释同样的 Harnack 因子。

\ProseParagraphBreak

\term[caratheodory-xishu-budengshi]{Carathéodory 系数不等式}[Carathéodory coefficient inequality]进一步给出与阶数无关的系数界。

\Needspace{6\baselineskip}
\begin{proposition}[Carathéodory 系数估计]\label{b3:prop:06-caratheodory-coefficients}
设 \(p:\mathbb D\rightarrow\mathbb C\) 全纯，\(\operatorname{Re}p>0\)、\(p(0)=1\)，且
\[
 p(z)=1+\sum_{n=1}^{\infty}c_nz^n.
\]
则每个整数 \(n\ge1\) 都满足 \(|c_n|\le2\)。
\end{proposition}
\begin{proof}
固定 \(0<r<1\)，记 \(u(t)=\operatorname{Re}p(re^{\mathrm{i}t})>0\)。Taylor 级数在这一固定圆周上一致收敛，故可取实部，写为
\[
 u(t)=1+\frac12\sum_{m=1}^{\infty}
 \left(c_mr^me^{\mathrm{i}mt}+\overline{c_m}r^me^{-\mathrm{i}mt}\right).
\]
对这个一致收敛级数逐项积分，利用非零整数频率的指数在一周上的积分为零，得到
\[
 \frac1{2\pi}\int_0^{2\pi}u(t)\,\mathrm{d}t=1,
 \quad
 \frac1{2\pi}\int_0^{2\pi}u(t)\cdot e^{-\mathrm{i}nt}\,\mathrm{d}t
 =\frac12c_nr^n\quad(n\ge1).
\]
因为 \(u(t)>0\) 且 \(|e^{-\mathrm{i}nt}|=1\)，带振荡因子的积分模不超过 \(u\) 的积分。因此 \(|c_n|\cdot r^n\le2\)。对固定 \(n\) 令 \(r\uparrow1\)，得到 \(|c_n|\le2\)。整个论证只在内部圆周上积分，不需要 \(p\) 在单位圆上具有边界值。
\end{proof}

上面的极端函数还有展开
\[
 \frac{1+\eta z}{1-\eta z}
 =1+2\sum_{n=1}^{\infty}\eta^nz^n\quad(|\eta|=1),
\]
所有非恒定系数的模均为 \(2\)，所以每个阶数都达到同一界。若全纯函数 \(p(z)=c_0+\sum_{n\ge1}c_nz^n\) 仍满足 \(\operatorname{Re}p>0\)，但常数项任意，令 \(\alpha=\operatorname{Re}c_0>0\)、\(\beta=\operatorname{Im}c_0\)。归一化函数 \((p-\mathrm{i}\beta)/\alpha\) 满足前一命题，故一般形式为 \(|c_n|\le2\operatorname{Re}c_0\)、\(n\ge1\)。常数项的虚部不参与这个约束。

\subsection{共形映射下的转移}
\label{b3:subsec:060304}

对真单连通区域，Riemann 映照提供全局圆盘坐标。先把距离和局部尺度转回原区域，并证明它们不依赖坐标选择，就能将 Schwarz--Pick 写成区域本身的结论。

\Needspace{6\baselineskip}
\begin{definition}[区域上的距离与密度]\label{b3:def:06-conformal-transfer}
设 \(U\subsetneq\mathbb C\) 为单连通区域，\(\Phi:U\rightarrow\mathbb D\) 为双全纯映射。对 \(z,w\in U\)，定义
\[
 \delta_U(z,w)\coloneqq\delta(\Phi(z),\Phi(w)),
 \quad
 \lambda_U(z)\coloneqq\frac{2|\Phi'(z)|}{1-|\Phi(z)|^2}.
\]
分别称为 \(U\) 上的伪双曲距离和\term[poincare-midu]{Poincaré 密度}[Poincaré density]。
\end{definition}
\symidx{\(\delta_U(z,w)\)：真单连通区域上的伪双曲距离}
\symidx{\(\lambda_U(z)\)：Poincaré 密度}

这两个定义与 \(\Phi\) 的选择无关。若另一坐标为 \(\Psi\)，则 \(A=\Psi\circ\Phi^{-1}\) 是圆盘自同构，前节的不变性给出 \(\delta(Au,Av)=\delta(u,v)\)。同时，自同构的导数恒等式为
\[
 \frac{|A'(u)|}{1-|A(u)|^2}=\frac1{1-|u|^2}.
\]
代入 \(u=\Phi(z)\) 并使用链式法则，便得用 \(\Psi\) 和用 \(\Phi\) 计算的 \(\lambda_U\) 相同。距离的度量性质由双射 \(\Phi\) 和 \(\delta\) 的性质传回；密度则处处为正，因为双全纯映射的导数不为零。

\ProseParagraphBreak

密度与上一章的共形半径可以在基点处直接比较。对固定 \(a\in U\)，选择满足 \(\Phi(a)=0\)、\(\Phi'(a)>0\) 的归一化坐标，得到
\[
 \lambda_U(a)=2\Phi'(a)=\frac2{r_U(a)}.
\]
坐标无关性保证这个等式对每个基点都成立。

\Needspace{6\baselineskip}
\begin{proposition}[单连通区域上的 Schwarz--Pick 估计]\label{b3:prop:06-conformal-schwarz-pick}
设 \(U,V\subsetneq\mathbb C\) 为单连通区域，\(f:U\rightarrow V\) 全纯。则对所有 \(z,w\in U\)，有
\[
 \delta_V(f(z),f(w))\le\delta_U(z,w),
 \quad \lambda_V(f(z))\cdot|f'(z)|\le\lambda_U(z).
\]
若在一对不同点的距离式中取等，或在某一点的导数式中取等，则 \(f\) 双全纯映满 \(V\)。反过来，双全纯映射在这些式子中处处取等。
\end{proposition}
\begin{proof}
取圆盘坐标 \(\Phi_U:U\rightarrow\mathbb D\)、\(\Phi_V:V\rightarrow\mathbb D\)，并令 \(G=\Phi_V\circ f\circ\Phi_U^{-1}\)。圆盘的 Schwarz--Pick 距离式直接变成所需的距离式。对 \(a\in U\)，在 \(\zeta=\Phi_U(a)\) 处有
\[
 G'(\zeta)=\frac{\Phi_V'(f(a))\cdot f'(a)}{\Phi_U'(a)}.
\]
将它代入圆盘导数式，再乘以正数 \(2|\Phi_U'(a)|\)，便得到 \(\lambda_V(f(a))\cdot|f'(a)|\le\lambda_U(a)\)。这里坐标导数均非零，所有分母为正，所以等号也准确对应。圆盘中的等号迫使 \(G\) 为自同构，故 \(f=\Phi_V^{-1}\circ G\circ\Phi_U\) 双全纯映满 \(V\)。反向同样由圆盘自同构的等号性质得到。
\end{proof}

用共形半径改写导数式，就是
\[
 \frac{|f'(a)|}{r_V(f(a))}\le\frac1{r_U(a)}.
\]
若 \(U\subseteq V\)，对恒等嵌入应用此式，得到 \(r_U(a)\le r_V(a)\)；一处取等就迫使恒等嵌入满射，所以只能有 \(U=V\)。这也恢复了上一章共形半径的严格单调性。

\ProseParagraphBreak

在右半平面 \(H_R=\{\,z\in\mathbb C\mid\operatorname{Re}z>0\,\}\)，取 \(\Phi(z)=(z-1)/(z+1)\)。由
\[
 |\Phi'(z)|=\frac2{|z+1|^2},
 \quad 1-|\Phi(z)|^2=\frac{4\operatorname{Re}z}{|z+1|^2},
\]
可得 \(\lambda_{H_R}(z)=1/\operatorname{Re}z\)。因此每个全纯映射 \(F:\mathbb D\rightarrow H_R\) 满足
\[
 \frac{|F'(z)|}{\operatorname{Re}F(z)}\le\frac2{1-|z|^2}.
\]
若 \(F(0)=1\)，距离式还给 \(|(F(z)-1)/(F(z)+1)|\le|z|\)，正是前面正实部估计所使用的圆盘约束。

\ProseParagraphBreak

在上半平面 \(\mathbb H=\{\,z\in\mathbb C\mid\operatorname{Im}z>0\,\}\)，同样用 \((z-\mathrm{i})/(z+\mathrm{i})\) 计算，得到
\[
 \delta_{\mathbb H}(z,w)=\left|\frac{z-w}{z-\overline w}\right|,
 \quad \lambda_{\mathbb H}(z)=\frac1{\operatorname{Im}z}.
\]
于是全纯映射 \(F:\mathbb H\rightarrow\mathbb H\) 满足
\[
 \left|\frac{F(z)-F(w)}{F(z)-\overline{F(w)}}\right|
 \le\left|\frac{z-w}{z-\overline w}\right|,
 \quad \frac{|F'(z)|}{\operatorname{Im}F(z)}\le\frac1{\operatorname{Im}z}.
\]
例如 \(F(z)=z+\mathrm{i}\) 的欧氏导数为 \(1\)，但 \(\operatorname{Im}F(z)=\operatorname{Im}z+1\)，所以密度导数式严格压缩。若 \(F(z)=(az+b)/(cz+d)\)，其中 \(a,b,c,d\) 为实数、\(ad-bc>0\)，则
\[
 \operatorname{Im}F(z)=\frac{(ad-bc)\operatorname{Im}z}{|cz+d|^2}.
\]
它及其分式逆都保持上半平面，故为自同构，在两个估计中均取等。

\ProseParagraphBreak

本小节的全局圆盘坐标依赖真单连通条件。对有孔区域，局部坐标不能直接拼成一个双全纯圆盘映射，推广密度需要后面的覆盖理论。\chapref{b3:ch:13}还将用 \(\lambda_U(z)\cdot|\mathrm{d}z|\) 定义曲线长度，再由长度下确界构造 Poincaré 距离；这里的 \(\delta_U\) 与那个路径距离是不同的数值。
